\begin{abstract}
Pooling couples network flows by requiring every outlet of a pool to
carry the same mixture. We prove that rational pooling threshold
decision is complete for the existential theory of the reals, both
with one attribute and no bypasses and with one pool and a growing
attribute dimension. The reductions preserve small degrees and bounded
data and imply that feasible flows can require unbounded algebraic
degree. With a fixed number of mixing parameters, basis indices instead
give polynomial-size certificates verified by real-algebraic decision.
We then establish matched hardness and tractability boundaries. Strong
NP-completeness persists with every local degree equal to two when the
pool count grows. With one pool, two feeds and two outlets, degree-three
products permit strong hardness with fixed finite physical data, whereas
degree-two products admit polynomial feasibility under the stated input
degree bound; the degree-three contracted version is also strongly
NP-complete for feasibility. Exact algorithms follow from independent bounded blocks,
small path boundaries, or conservation under exact contracts. These
algorithms specify their objective scope and algebraic output. Supporting
constructions separate representation complexity from optimization:
physical paths have exponentially many price responses, and a one-pool
family has exponentially many isolated optima yet a compact exact LP
hull. For an isolated rank-one flow block, we resolve the general-cost
hardness conjecture, prove exact tractability for fixed interaction rank,
and establish exact and uniform approximate conic size lower bounds.
\end{abstract}

\section*{Introduction}
\addcontentsline{toc}{section}{Introduction}
A pool combines incoming streams and sends the same mixture to every
outlet. Once pool compositions are fixed, flow conservation, capacities,
and product quality specifications form a linear system. Choosing the
compositions introduces bilinear equations. Direct input-to-product
arcs, called bypasses, further couple pool intakes with material blended
directly at products through shared source capacities. These interactions
make pooling a basic model of nonconvex network optimization.

The established complexity literature identifies NP-hard restrictions
and polynomial special cases. It leaves a more basic issue: does exact
pooling decision have short discrete certificates at all? A rational
input need not have a rational feasible flow. Our principal result places
pooling threshold decision in the class $\ETR$ of problems reducible to
the existential theory of the reals, and proves completeness under two
restrictions: one attribute without bypasses, or one pool with bypasses
and growing attribute dimension. The first construction also admits
bounded degrees and fixed finite data refinements. The reductions give
rational correspondences between solution sets and transfer algebraic
witness obstructions. Thus general threshold decision belongs to $\NP$
only if $\NP=\ETR$; irrationality alone would not imply that conclusion.

\citet{haugland2016-the-computational-complexity-of-the} established
NP-hardness under restrictions on network cardinalities. Our algebraic
reductions address a different barrier, using exact product quality
specifications to realize addition and inversion. To the best of our
knowledge, the pooling literature compared in Section~\ref{s2:algebraic}
contains no previous $\ETR$-completeness result for the two stated
classes. When pool count times affine input-quality rank, or pool count
times product count, is fixed, a guessed LP basis instead supplies a
polynomial-size certificate verified by fixed-dimensional algebraic
decision. This is $\NP$ membership, not a polynomial optimization
algorithm or a claim of rational optimal flows.

Our second contribution identifies combinatorial hardness under
restrictions close to the algorithmic regimes. Boland, Kalinowski and
Rigterink \cite[author manuscript, Section~4]{s4:boland2017} ask whether
scalar-quality pooling becomes polynomial when all in-degrees, or all
out-degrees, are at most two. We prove strong NP-completeness when all
four layer degrees are exactly two simultaneously; the profit problem
also has no PTAS unless $\PP=\NP$. Their separate bounded-two-pool
question, also raised by Haugland, is answered negatively with two
products on an undirected tree, allowing input and attribute counts to
grow. With bypasses, we prove strong hardness with one upper-bounded
scalar quality, two actual feeds, two actual outlets, and fixed finite
physical data. Broad one-pool capacity hardness was already stated by
\citet[Remark~4.6]{s3:baltean2018}; our advance is the explicit reduction
with these simultaneous restrictions. Many external bypass sources do
not count as additional mixing feeds.

The third contribution gives exact algorithms where the linear network
can be eliminated without losing pool balances or the objective.
A fixed nonlinear core coupled to independent polyhedral blocks of
fixed local dimension admits polynomial global optimization. Its pooling
consequences include fixed input and pool counts, and fixed pool count,
affine quality rank and bypass vertex integrity. For long degree-two
bypass paths, planar composition yields polynomial endpoint relations
when the retained boundary is fixed. In particular, with one pool, two
feeds, two outlets, and input degree at most two, product degree two
admits polynomial feasibility with arbitrary attributes and positive
flow lower bounds. With product degree three, the contracted construction is strongly
NP-complete for feasibility, with the same two-feed, two-outlet
interface and input degree bound (\cref{s4:degree-boundary}). This
is the matched feasibility boundary. The zero-flow-lower-bound
economic threshold result of \cref{s3:constant-thm} is a separate
refinement, and the degree-two algorithm does not cover dense objectives.

Exact contracts provide a different way to retain global information.
They turn pool mass and attribute mass into conserved boundary quantities,
or into signed node divergences on degree-two bypass graphs. We obtain
polynomial algorithms with a fixed number of receiving products, or a
fixed number of exceptions to otherwise exact external contracts.
The latter allows arbitrary quality rank when common pool bounds are
redundant. Fully contracted scalar or affine-rank-one instances admit
restrictive common throughput bounds and linear throughput objectives.
With only two distinct full input-quality vectors, a separate convex
reduction gives rational feasibility witnesses on arbitrary bypass
graphs under zero outlet and common lower bounds. These cases retain
hypotheses absent from general hardness: fixed pool count, fixed rank,
small retained boundaries, or exact conservation. The detailed scope
comparison in Table~\ref{s5:scope-table} prevents combining restrictions
from different theorems without an additional argument.

Table~\ref{s0:boundaries} gives the main physical boundaries. A fixed
parameter is an input-independent constant; polynomial exponents may
depend on it. No fixed-parameter tractability claim is implicit.
\emph{Standard economics} means linear source production costs and
product revenues. Other arc costs are allowed only where stated.
An exact product contract fixes demand and the full quality vector;
at zero flow the quality equations have their homogeneous meaning.

\begin{table}[!t]
\centering
\small
\setlength{\tabcolsep}{4pt}
\begin{tabular}{@{}>{\raggedright\arraybackslash}p{.39\linewidth}>{\raggedright\arraybackslash}p{.36\linewidth}>{\raggedright\arraybackslash}p{.19\linewidth}@{}}
\toprule
Physical restriction & Decision or objective scope & Result\\
\midrule
One attribute, no bypasses; or one pool and growing attribute dimension
& Rational thresholds, two-sided qualities, zero flow lower bounds
& $\ETR$-complete (\cref{s2:etr-complete,s2:one-pool-etr})\\
\addlinespace
Fixed pool count times quality rank or product count
& Feasibility and linear arc-cost thresholds; unrestricted bypasses
& In $\NP$ (\cref{s2:pooling-np})\\
\addlinespace
All four layer degrees two; growing pool count
& Upper-only scalar qualities, zero flow lower bounds; profit
& Strongly $\NP$-complete; no PTAS (\cref{s3:all-two-thm})\\
\addlinespace
One pool, two feeds and two outlets, input degree two
& Feasibility with contracts allowed: arbitrary attributes at product
degree two; fixed scalar data suffice at degree three
& Degree two: polynomial; degree three: strongly $\NP$-complete
(\cref{s4:degree-boundary})\\
\addlinespace
Fixed pools, quality rank and bypass vertex integrity
& Feasibility and arbitrary linear arc-cost optimization
& Polynomial (\cref{s4:vertex-integrity})\\
\addlinespace
Degree-two bypasses; fixed total pool attachments
& Feasibility and objectives on a fixed retained arc set
& Polynomial (\cref{s4:bounded-attachments,s4:fixed-support})\\
\addlinespace
One pool with exact contracts, bounded exceptions, or two full input vectors
& Distinct feasibility and objective scopes in Table~\ref{s5:scope-table}
& Polynomial under the listed hypotheses\\
\bottomrule
\end{tabular}
\caption{Selected physical boundaries, with finite rational capacities.
Hardness and algorithmic assumptions are completed in the cited statements.
The no-PTAS conclusion assumes $\PP\ne\NP$.}
\label{s0:boundaries}
\end{table}

Two secondary contributions clarify what a complexity classification
can and cannot imply. Physical blending paths can have exponentially
many responses to one varying product price. A one-pool extension has
exponentially many isolated global and strict local optima, yet a
polynomial-size exact LP hull and polynomial linear optimization.
These constructions separate explicit response descriptions and local
geometry from computational hardness. For the isolated rank-one margin
block used in pooling relaxations, we prove the general-cost hardness
conjecture of \citet{dey2020-convexifications-of-rank-one-based}, together
with polynomial exact optimization for fixed cost interaction rank.
We also embed a correlation-polytope face and prove quantitative
transfers giving exact and uniform approximate conic size lower bounds.
These statements concern arbitrary matrix costs or a uniform hull;
ordinary additive operating costs and a single chosen objective have
different scope. The exact exponential SOC representation of
\citet{s6:jalilian2026} is consistent with all these bounds.

The paper also supplies physical proofs of useful established tools.
Xiang et al.\ \cite{xiang2026-dive-and-fix} already give a polynomial
zero-optimum test by one LP per product, including pool-to-pool arcs.
Section~\ref{s1:foundations} credits that result and develops its
physical decomposition, conic-hull interpretation, and explicit
circulation conventions.
Earlier polynomial pooling cases include
\citet{s3:haugland-hendrix2016,s4:boland2017}; their cardinality and
bypass assumptions are compared where used. The parametric analysis of
\citet{s3:baltean2018} assumes unbounded feed availability and pool
capacity and fixed positive demands. Approximation and formulation work
by \citet{dey2015-analysis-of-milp-techniques-for,
gupte2017-relaxations-and-discretizations-for-the,
dey2020-convexifications-of-rank-one-based} supplies further context.

Section~\ref{s1:foundations} defines the models and exact arithmetic.
Sections~\ref{s2:algebraic} and~\ref{s3:hardness} prove algebraic
completeness and restricted hardness; Sections~\ref{s4:structural}
and~\ref{s5:contracts} develop structural and contract algorithms.
Section~\ref{s6:synthesis} explains their representation consequences
and the remaining degree-two bypass question. Appendices~\ref{app:geometry},
\ref{app:costs} and~\ref{app:convexification} give complete proofs for
physical response geometry, isolated-block costs, and convexification.

